% =====================================================================
\section{Training Procedure and Experiment Management}
\label{sec:training}
% =====================================================================

\subsection{Training workflow}
Training is driven by a single script, \texttt{train.py}, that reads one YAML
configuration file and performs the following steps.

\begin{enumerate}
  \item \textbf{Set up.} Load the configuration, set all random seeds, create a
        run directory named after the experiment and timestamp, and copy the
        configuration and Git commit hash into it.
  \item \textbf{Build datasets.} Filter \texttt{manifest.csv} by the speakers
        listed for each partition in \texttt{splits.csv}. The training dataset
        returns, per utterance, the cached mouth crops, the clean waveform mixed
        with noise at a randomly drawn SNR, and the encoded transcript. Validation
        and test datasets read the fixed pre-mixed noisy files.
  \item \textbf{Build batches.} A custom collate function pads waveforms and crop
        sequences, computes feature and encoder lengths, builds padding masks and
        concatenates targets for CTC (\cref{sec:audio}). Log-Mel extraction and
        normalisation run on the GPU on the padded batch.
  \item \textbf{Initialise the model.} Convolution and linear layers use PyTorch's
        default (Kaiming-uniform) initialisation; LSTM forget-gate biases are set
        to 1 to ease gradient flow early in training; the CTC output bias for the
        blank symbol is left at zero.
  \item \textbf{Train.} For each batch: forward pass (\cref{alg:forward}), CTC loss
        (\cref{eq:ctc}) averaged over target lengths, backward pass with automatic
        mixed precision (AMP), gradient-norm clipping, optimiser step and
        scheduler step.
  \item \textbf{Validate.} After each epoch, compute validation CTC loss and
        greedy-decoded WER and CER at every SNR, and their mean.
  \item \textbf{Checkpoint.} Save the latest state and, separately, the best state
        by mean validation WER. A checkpoint contains model, optimiser,
        scheduler and AMP scaler states, the epoch, the configuration and the
        vocabulary.
  \item \textbf{Stop.} End training when mean validation WER has not improved for
        the patience period, or at the maximum epoch count.
  \item \textbf{Restore and export.} Reload the best checkpoint, run the final
        evaluation (\cref{sec:evaluation}) and export the model weights for the
        demonstration.
\end{enumerate}

\subsection{Initial hyperparameters}
\Cref{tab:hparams} lists starting values \rec. They are common defaults for
models of this size and are \emph{not} established optimal settings for
LipScribe; the most sensitive values (learning rate, dropout and modality-dropout
rates) will be tuned on the validation set.

\begin{table}[H]
\centering
\caption{Initial training hyperparameters (starting values to be tuned).}
\label{tab:hparams}
\small
\begin{tabularx}{\textwidth}{@{}p{4.2cm}p{3.4cm}X@{}}
\toprule
\textbf{Hyperparameter} & \textbf{Starting value} & \textbf{Rationale} \\
\midrule
Optimiser & AdamW & Robust default for recurrent and attention models \\
Peak learning rate & $3\times10^{-4}$ & Tuned over $\{10^{-4}, 3\times10^{-4}, 10^{-3}\}$ \\
Weight decay & $10^{-2}$ & Mild regularisation; not applied to biases or norms \\
Schedule & 1 epoch linear warm-up, then cosine decay & Warm-up stabilises early CTC training \\
Batch size & 16 utterances & Expected to fit an 8\,GB GPU with AMP; use gradient accumulation otherwise \\
Gradient clipping & global norm 5.0 & Prevents LSTM gradient explosions \\
Dropout & 0.3 (encoders), 0.1 (attention) & Limits overfitting to 27 training speakers \\
Modality dropout & $p_a=0.25$, $p_v=0.15$ & \Cref{sec:modality-dropout} \\
Maximum epochs & 60 & Upper bound; early stopping usually ends sooner \\
Early-stopping patience & 8 epochs & Measured on mean validation WER \\
Beam width (evaluation) & 10 & \Cref{sec:ctc} \\
Random seeds & 1, 2, 3 & Three seeds for the final models, if compute allows \\
\bottomrule
\end{tabularx}
\end{table}

\paragraph{Gradient management.} LSTMs and CTC can produce occasional very large
gradients, particularly early in training. Global-norm clipping bounds the update
size. With AMP, gradients are unscaled before clipping. If a batch produces a
non-finite loss, the step is skipped and logged; more than a handful of such
events indicates a data or length bug rather than bad luck.

\paragraph{Curriculum for visual models.} Visual-only CTC training is known to be
slow to start. If the visual-only baseline has not begun to reduce validation
CER after several epochs, the team will first train the visual encoder alone and
then use its weights to initialise the visual stream of the fused model \tbd.

\subsection{Experiment logging and management}
Each run logs training loss, validation loss, WER and CER per SNR, learning rate,
gradient norm and mean gate value to TensorBoard, and appends one summary row to
\texttt{experiments/results.csv} on completion. Run directories are never
overwritten. An experiment is reported only if its configuration file is
committed, so that every number in the final report can be traced to a
configuration, a commit and a checkpoint.

\begin{lstlisting}[language={},caption={Abbreviated experiment configuration for the proposed model \rec.},label={lst:config}]
experiment: proposed_attention_fusion
seed: 1
data:
  manifest: data/manifest.csv
  splits: data/splits.csv
  train_snrs: [clean, 10, 5, 0, -5]
  eval_snrs:  [clean, 10, 5, 0, -5]
model:
  streams: [audio, visual]
  visual_encoder: conv3d_bilstm      # V1; alternative: conv3d_effnetb0
  audio_features: logmel80           # A1; alternatives: mfcc39, logmel_mfcc
  fusion: cross_attention_gated      # alternatives: concat, none
  modality_dropout: {audio: 0.25, visual: 0.15}
train:
  optimiser: adamw
  lr: 3.0e-4
  batch_size: 16
  max_epochs: 60
  patience: 8
  grad_clip: 5.0
decode: {method: beam, beam_width: 10}
\end{lstlisting}

\subsection{Baselines and experiment plan}
The proposal requires three simple baselines---audio-only, visual-only and
simple concatenation---followed by the full system with a 3D convolutional
visual encoder and adaptive cross-modal attention with modality
dropout~\spec. All models share the same code, encoders, data partitions,
noise files, feature normalisation, optimiser, schedule, maximum epoch count and
early-stopping rule. Only the components under test differ. The audio-only
baseline is trained with the same multi-condition noise as the fused models; a
baseline trained only on clean audio would be unrealistically weak in noise and
make the comparison unfair.

\begin{itemize}
  \item \textbf{B1 Audio-only:} audio encoder $\rightarrow$ sequence BiLSTM $\rightarrow$ CTC.
  \item \textbf{B2 Visual-only:} visual encoder (V1) $\rightarrow$ sequence BiLSTM $\rightarrow$ CTC.
  \item \textbf{B3 Simple fusion:} concatenate $\mathbf{A}$ and $\mathbf{V}$ per frame,
        linear projection to 256, then the same sequence BiLSTM and CTC; no attention,
        no modality dropout.
  \item \textbf{P Proposed:} the complete model of \cref{sec:model}.
\end{itemize}

\begin{table}[H]
\centering
\caption{Planned experiments. No results are reported; the last column states what
each comparison is designed to reveal.}
\label{tab:experiments}
\small
\begin{tabularx}{\textwidth}{@{}p{0.6cm}p{3.2cm}p{3.6cm}X@{}}
\toprule
\textbf{ID} & \textbf{Model} & \textbf{Key comparison} & \textbf{Diagnostic value} \\
\midrule
B1 & Audio-only & Reference for all fused models & How fast ASR alone degrades as SNR falls \\
B2 & Visual-only & Constant across SNR by design & Upper bound on what the lips alone can provide on GRID; checks the visual pipeline works \\
B3 & Concatenation fusion & B3 vs.\ B1 & Whether adding video helps at all without adaptive weighting \\
P  & Attention fusion + modality dropout & P vs.\ B3, P vs.\ B1 & Benefit of adaptive fusion, especially at 0 and $-5$\,dB \\
X1 & P without modality dropout & X1 vs.\ P & Contribution of modality dropout; whether video is otherwise ignored \\
X2 & P with concatenation in place of attention, keeping dropout & X2 vs.\ P & Contribution of attention itself, separated from dropout \\
X3 & P with a 2D per-frame CNN instead of the 3D front-end & X3 vs.\ P & Contribution of 3D (spatiotemporal) convolution \\
X4 & P with EfficientNet-B0 trunk (V2) & X4 vs.\ P & Value of a pretrained 2D backbone (optional) \\
X5 & B1 with A2 / A3 features & vs.\ B1 & Whether MFCC features add anything to log-Mel \\
X6 & P with SNR auxiliary head & X6 vs.\ P & Whether an explicit noise estimate improves gating (optional) \\
\bottomrule
\end{tabularx}
\end{table}

B1, B2, B3 and P form the core and are completed first. X1--X3 are the ablations
the proposal asks for and are run next. X4--X6 are optional and depend on
remaining time and compute.
